\subsection{Regular extensions}

DEFINITION 2. --- An extension of a field K is said to be regular if it is regular as K-algebra.

PROPOSITION 4. --- Let A be an algebra over field K which is an integral domain, and E its field of fractions. Let L be an extension of K; if the ring \(L \otimes_{K} A\) is an integral domain the same is true of \(L \otimes_{K} E\) .

If \(L \otimes_{K} A\) is an integral domain, it may be embedded in its field of fractions F. Write \(u(x) = x \otimes 1\) for \(x \in L\) and denote by v the K-homomorphism of E into F which extends the injective homomorphism \(y \mapsto 1 \otimes y\) of A into F. By Prop. 6 (V, p.~14) the subfields \(u(L)\) and \(v(E)\) of F are linearly disjoint over K; therefore the homomorphism \(u * v\) of \(L \otimes_{K} E\) into F (V, p.~12) is injective. This shows \(L \otimes_{K} E\) to be an integral domain.

COROLLARY. --- For A to be a regular K-algebra it is necessary and sufficient that its field of fractions should be a regular extension of K.

The condition is necessary by Prop. 4 and sufficient by Prop. 3, a).

PROPOSITION 5. --- Every pure extension of a field K is regular.

By the preceding corollary it is enough to prove that every polynomial algebra \(A = K[X_{i}]_{i \in I}\) is a regular K-algebra. Let L be an extension of \(K\) ; the ring \(L \otimes_{K} A\) is isomorphic to \(L[X_{i}]_{i \in I}\) (III, p.~449, Remark 2), and hence an integral domain (IV, p.~9, Prop. 8).

PROPOSITION 6. --- Let L be an extension of a field K. If L is regular, then every subextension of L is regular. Conversely, if every finitely generated subextension of L is regular, then L is regular.

The first assertion follows from Prop. 3, a).

Let M be an extension of K and let U be the set of all finitely generated subextensions of L. For each \(E \in U\) , the ring \(M \otimes_{K} E\) may be identified with a subring of \(M \otimes_{K} L\) and we thus have an increasing directed family of subrings of \(M \otimes_{K} L\) whose union is \(M \otimes_{K} L\) . Now the second assertion follows immediately.

PROPOSITION 7. --- Let L be an extension of a field K and M an L-algebra (for example, an extension of L). If L is regular over K and M is a regular L-algebra, then M is regular as K-algebra.

Let E be an extension of K; since L is regular over K, \(E \otimes_{K} L\) is an integral domain. By Prop. 2 (V, p.~141) the ring \((\mathrm{E} \otimes_{\mathrm{K}} \mathrm{L}) \otimes_{\mathrm{L}} \mathrm{M}\) is thus an integral domain and the same is true of the ring \(E \otimes_{K} M\) isomorphic to it (II, p.~278, Prop. 2). Hence the result.

PROPOSITION 8. --- Let L and M be two extensions of a field K.

\begin{enumerate}
\def\labelenumi{\alph{enumi})}
\item
  If \(\mathbf{M}\) is regular over \(\mathbf{K}\) , then the field of fractions of the integral domain \(\mathbf{L} \otimes_{\mathbf{K}} \mathbf{M}\) is a regular extension of \(\mathbf{L}\) .
\item
  If L and M are regular extensions of K, the same is true of the field of fractions of L ⊗\_K M.
\end{enumerate}

Assertion a) follows from Prop. 3, c) (V, p.~141) and the Cor. of V, p.~141; Assertion b) follows from Prop. 3, b) (V, p.~141) and the Cor. of V, p.~141.

